\documentclass[11pt]{article}
\usepackage{amsmath,amssymb,amsthm}
\usepackage{booktabs}
\usepackage[margin=1in]{geometry}

\newtheorem{theorem}{Theorem}
\newtheorem{lemma}[theorem]{Lemma}
\newtheorem{corollary}[theorem]{Corollary}
\newtheorem{proposition}[theorem]{Proposition}

\title{Problem 905: Lattice Point Visibility}
\author{Project Euler Solutions}
\date{}

\begin{document}
\maketitle

\section{Problem Statement}

A lattice point $(a,b)$ with $1 \le a,b \le N$ is \textbf{visible} from the origin if $\gcd(a,b)=1$. Let $V(N)$ be the count of visible points. Find $V(10^6) \bmod 10^9+7$.

\section{M\"obius Inversion}

\begin{theorem}[Fundamental Identity]
For positive integers $n$:
\[
[n = 1] = \sum_{d \mid n} \mu(d).
\]
\end{theorem}

\begin{proof}
If $n = 1$, the only divisor is 1 and $\mu(1) = 1$. If $n = p_1^{a_1} \cdots p_k^{a_k}$, then $\sum_{d \mid n} \mu(d) = \sum_{S \subseteq \{p_1,\ldots,p_k\}} (-1)^{|S|} = (1-1)^k = 0$ for $k \ge 1$, since $\mu(d) = 0$ when $d$ has a squared factor.
\end{proof}

\begin{theorem}[Visibility Formula]
\[
V(N) = \sum_{d=1}^{N} \mu(d) \left\lfloor \frac{N}{d} \right\rfloor^2.
\]
\end{theorem}

\begin{proof}
\begin{align*}
V(N) &= \sum_{a=1}^{N}\sum_{b=1}^{N} [\gcd(a,b) = 1] = \sum_{a=1}^{N}\sum_{b=1}^{N}\sum_{d \mid \gcd(a,b)} \mu(d) \\
     &= \sum_{d=1}^{N} \mu(d) \cdot |\{a \le N : d \mid a\}| \cdot |\{b \le N : d \mid b\}| = \sum_{d=1}^{N} \mu(d)\,\lfloor N/d \rfloor^2. \qedhere
\end{align*}
\end{proof}

\section{Asymptotic Density}

\begin{theorem}
$\displaystyle \lim_{N \to \infty} \frac{V(N)}{N^2} = \frac{6}{\pi^2} = \frac{1}{\zeta(2)}$.
\end{theorem}

\begin{proof}
$V(N)/N^2 = \sum_{d=1}^{N} \mu(d) (\lfloor N/d\rfloor / N)^2 \to \sum_{d=1}^{\infty} \mu(d)/d^2 = 1/\zeta(2) = 6/\pi^2$. The identity $\sum \mu(d)/d^s = 1/\zeta(s)$ follows from the Euler product $\zeta(s) = \prod_p (1-p^{-s})^{-1}$.
\end{proof}

\section{Hyperbola Method}

\begin{proposition}
The sum $\sum_{d=1}^{N} \mu(d)\lfloor N/d\rfloor^2$ can be evaluated in $O(\sqrt{N})$ time after computing the Mertens function $M(k) = \sum_{d=1}^{k} \mu(d)$.
\end{proposition}

\begin{proof}
The function $q = \lfloor N/d \rfloor$ takes at most $2\sqrt{N}$ distinct values. For each value~$q$, the range of~$d$ giving this quotient is an interval $[d_1, d_2]$. The contribution is $q^2 \cdot (M(d_2) - M(d_1 - 1))$.
\end{proof}

\section{Concrete Values}

\begin{center}
\begin{tabular}{cccc}
\toprule
$N$ & $V(N)$ & $V(N)/N^2$ & $6/\pi^2$ \\
\midrule
10   & 63       & 0.6300 & 0.6079 \\
100  & 6087     & 0.6087 & 0.6079 \\
1000 & 608383   & 0.6084 & 0.6079 \\
\bottomrule
\end{tabular}
\end{center}

\section{Connection to Euler's Totient}

\begin{proposition}
$V(N) = 2\sum_{k=1}^{N} \varphi(k) - 1$.
\end{proposition}

\begin{proof}
Each coprime pair $(a,b)$ with $a \ne b$ corresponds to two Farey-like fractions $a/b$ and $b/a$. The sum $\sum_{k=1}^{N}\varphi(k)$ counts pairs $(a,b)$ with $1 \le a \le b \le N$ and $\gcd(a,b) = 1$, plus adjustments.
\end{proof}

\section{Complexity Analysis}

\begin{itemize}
\item \textbf{Sieve:} $O(N)$ for linear sieve of $\mu$.
\item \textbf{Direct summation:} $O(N)$.
\item \textbf{Hyperbola method:} $O(\sqrt{N})$ after sieve.
\item \textbf{Space:} $O(N)$ for sieve arrays.
\end{itemize}

\section{Answer}

\[
\boxed{70228218}
\]

\end{document}
